% ======================================================================
% Appendix G. Fourier normalizations and spectral measures
% PATH: src/appendices/G-fourier-normalizations-and-measures.tex
% ======================================================================

\appendix
\chapter{Fourier normalizations, spectral measures, and compatibility conventions}
\label{app:G.fourier}

\section{Scope and guiding principle}
\label{sec:G.scope}

The purpose of this appendix is to fix, once and for all, the Fourier transform conventions and the associated
spectral measures used throughout the monograph.  Since localized trace identities are highly sensitive to
normalization choices, even a minor mismatch in factors of $2\pi$ or in the parametrization of the spectrum can
propagate into erroneous leading terms.  The role of the present appendix is therefore structural: it provides a
single authoritative reference against which all Fourier-analytic expressions in the main text are synchronized.

The conventions recorded below are standard, but we insist on explicitly stating them, because the main results are
formulated at a quantitative level where the leading constants and power-saving remainders must remain invariant under
admissible modifications.

\section{Fourier transform on $\R$}
\label{sec:G.fourier.R}

\subsection{Definition}
\label{subsec:G.fourier.def}

For $f\in \mathcal{S}(\R)$ we define the Fourier transform by
\begin{equation}\label{eq:G.fourier.def}
\widehat{f}(\xi)
\;:=\;
\int_{\R} f(t)\,e^{-i\xi t}\,dt,
\qquad
\xi\in\R.
\end{equation}
The inverse transform is
\begin{equation}\label{eq:G.fourier.inv}
f(t)
\;=\;
\frac{1}{2\pi}\int_{\R}\widehat{f}(\xi)\,e^{i\xi t}\,d\xi.
\end{equation}

This convention is used uniformly for all Fourier-analytic manipulations in the wave kernel representation,
stationary phase arguments, and Tauberian smoothing removal.

\subsection{Plancherel and convolution}
\label{subsec:G.fourier.plancherel}

With the convention \eqref{eq:G.fourier.def}--\eqref{eq:G.fourier.inv}, the Plancherel identity takes the form
\begin{equation}\label{eq:G.plancherel}
\int_{\R}|f(t)|^2\,dt
\;=\;
\frac{1}{2\pi}\int_{\R}|\widehat{f}(\xi)|^2\,d\xi.
\end{equation}

The convolution of two functions $f,g\in L^1(\R)$ is defined by
\begin{equation}\label{eq:G.convolution.def}
(f\ast g)(t)
\;:=\;
\int_{\R} f(t-s)\,g(s)\,ds,
\end{equation}
and satisfies the standard product rule
\begin{equation}\label{eq:G.convolution.rule}
\widehat{f\ast g}(\xi)=\widehat{f}(\xi)\,\widehat{g}(\xi).
\end{equation}

For $f^\ast(t):=\overline{f(-t)}$ we have
\begin{equation}\label{eq:G.convolution.square}
\widehat{f\ast f^\ast}(\xi)=|\widehat{f}(\xi)|^2\ge 0,
\end{equation}
a fact which underlies the quadratic-form positivity mechanisms used in Chapter~\ref{chap:10-trace-to-zeta}.

\subsection{Evenness and real-valuedness}
\label{subsec:G.fourier.even}

If $f$ is even, $f(t)=f(-t)$, then $\widehat{f}$ is even.  If, in addition, $f$ is real-valued, then $\widehat{f}$ is
real-valued and even.

These parity properties are invoked repeatedly in the construction of even test families $h(t)$ in the trace formula.

\section{Spectral parametrization for hyperbolic surfaces}
\label{sec:G.spectral.param}

Let $M=\Gamma\backslash\H$ be a finite-area hyperbolic surface with cusps, and let $\Delta$ denote the positive
Laplace--Beltrami operator.  The spectrum decomposes as
\[
\Spec(\Delta)
=
\Spec_{\mathrm{disc}}(\Delta)\;\cup\;[1/4,\infty),
\]
where the discrete eigenvalues are written in the form
\begin{equation}\label{eq:G.eigs}
\lambda_j=\frac14+r_j^2,
\qquad
r_j\in[0,\infty)\cup i(0,1/2].
\end{equation}

We use $r$ as the spectral parameter on the continuous spectrum, so that $\lambda=1/4+r^2$.
In particular, the spectral localization regimes in the main theorems are formulated in terms of large $|r|$.

\subsection{Spectral projector conventions}
\label{subsec:G.projector}

Let $h\in \mathcal{S}(\R)$ be even.  The functional calculus defines
\[
h(\sqrt{\Delta-1/4}),
\]
which acts diagonally on eigenfunctions and via Eisenstein series on the continuous spectrum.
The normalization adopted in the main text is consistent with the spectral decomposition
\begin{equation}\label{eq:G.spectral.decomp}
\Tr\big(h(\sqrt{\Delta-1/4})\big)
=
\sum_{j} h(r_j)
+
\frac{1}{4\pi}\sum_{\mathfrak{a}}
\int_{\R} h(r)\,\varphi_{\mathfrak{a}}(1/2+ir)\,dr,
\end{equation}
where $\varphi_{\mathfrak{a}}(s)$ denotes the logarithmic derivative of the scattering determinant at cusp
$\mathfrak{a}$ in the normalization fixed below.

The factor $1/(4\pi)$ is the standard Plancherel weight for Eisenstein series on $\Gamma\backslash\H$ and must be
regarded as part of the definition of the continuous spectral measure.

\section{Eisenstein series normalization and scattering data}
\label{sec:G.scattering}

\subsection{Eisenstein series and constant term}
\label{subsec:G.eisenstein}

Let $\mathfrak{a}$ be a cusp of $\Gamma$.  We denote by $E_{\mathfrak{a}}(z,s)$ the Eisenstein series normalized so that
its constant term at cusp $\mathfrak{b}$ takes the form
\begin{equation}\label{eq:G.constant.term}
E_{\mathfrak{a}}(z,s)\big|_{\mathfrak{b}}
=
\delta_{\mathfrak{a}\mathfrak{b}}\,y^s
+
\phi_{\mathfrak{a}\mathfrak{b}}(s)\,y^{1-s},
\end{equation}
where $y$ is the height function in the standard cusp coordinate.

The matrix $\Phi(s)=(\phi_{\mathfrak{a}\mathfrak{b}}(s))$ is the scattering matrix.
Its determinant is denoted by $\varphi(s)=\det\Phi(s)$.

\subsection{Logarithmic derivative convention}
\label{subsec:G.logder}

The function $\varphi_{\mathfrak{a}}(s)$ appearing in \eqref{eq:G.spectral.decomp} is defined by
\begin{equation}\label{eq:G.varphi.a}
\varphi_{\mathfrak{a}}(s)
\;:=\;
\frac{d}{ds}\log \varphi_{\mathfrak{a}}(s),
\end{equation}
where $\varphi_{\mathfrak{a}}(s)$ denotes the diagonal scattering coefficient (or, equivalently, the cuspwise
contribution to $\det\Phi(s)$) in the standard factorization.
In particular, for the total scattering determinant one has
\[
\frac{d}{ds}\log\varphi(s)
=
\sum_{\mathfrak{a}}\varphi_{\mathfrak{a}}(s).
\]

This is the convention used in Chapters~\ref{chap:06-parabolic} and~\ref{chap:07-main-results}.

\subsection{Functional equation}
\label{subsec:G.func.eq}

The scattering matrix satisfies the functional equation
\[
\Phi(s)\Phi(1-s)=I,
\]
and hence $\varphi(s)\varphi(1-s)=1$.
Differentiating yields
\begin{equation}\label{eq:G.logder.func.eq}
\frac{\varphi'}{\varphi}(s)
=
-\frac{\varphi'}{\varphi}(1-s).
\end{equation}
This symmetry is used repeatedly to reduce continuous-spectrum integrals to the critical line $\Re(s)=1/2$.

\section{The Selberg--Harish-Chandra transform}
\label{sec:G.selberg.transform}

\subsection{Radial kernels and their spectral transform}
\label{subsec:G.shc}

Let $k(z,w)=K(d(z,w))$ be a radial kernel on $\H$ depending only on hyperbolic distance.
Its Selberg--Harish-Chandra transform is the function $h(r)$ defined by
\begin{equation}\label{eq:G.shc.def}
h(r)
=
\int_0^\infty K(u)\,\varphi_r(u)\,(\sinh u)\,du,
\end{equation}
where $\varphi_r$ is the spherical function at spectral parameter $r$.

We do not fix a separate normalization for $\varphi_r$ here, since all spherical transforms in the main text are
expressed in terms of the standard Selberg transform conventions in \cite{Hejhal1983,Iwaniec2002}.
The present appendix serves to synchronize the $2\pi$-factors and Plancherel measures once the transform is chosen.

\subsection{Compatibility with the Euclidean Fourier transform}
\label{subsec:G.compatibility}

The Euclidean Fourier transform convention \eqref{eq:G.fourier.def} is compatible with the spherical transform in the
sense that the wave kernel representation (Chapter~\ref{chap:03-kernel}) may be written as
\begin{equation}\label{eq:G.wave.repr}
h(r)
=
\int_{\R}\widehat{h}(t)\,e^{itr}\,dt,
\end{equation}
with the same $\widehat{h}$ as in \eqref{eq:G.fourier.def}, provided $h$ is even.
This identity is interpreted distributionally for general $h$ and is made precise via cutoff approximations.

\section{Plancherel measure on $\Gamma\backslash\H$}
\label{sec:G.plancherel.hyperbolic}

\subsection{Local model: $\H$}
\label{subsec:G.plancherel.H}

On the hyperbolic plane $\H$, the spectral measure for the Laplacian corresponds to the density
\[
\frac{1}{4\pi}\,r\tanh(\pi r)\,dr,
\]
which is the standard Plancherel density for spherical analysis.

In localized trace identities, the factor $r\tanh(\pi r)$ is responsible for the leading Weyl term
$\asymp \vol(M)\,r$.

\subsection{Global quotient}
\label{subsec:G.plancherel.M}

On $M=\Gamma\backslash\H$, the discrete spectrum is counted with multiplicity, while the continuous spectrum carries
the measure appearing in \eqref{eq:G.spectral.decomp}.
In particular, for any even Schwartz function $h$ one has
\[
\Tr(h(\sqrt{\Delta-1/4}))
=
\sum_{j} h(r_j)
+
\frac{1}{4\pi}\sum_{\mathfrak{a}}
\int_{\R} h(r)\,\varphi_{\mathfrak{a}}(1/2+ir)\,dr,
\]
and this identity should be viewed as the definition of the global spectral measure compatible with our Fourier
conventions.

\section{Scaling conventions for localized windows}
\label{sec:G.scaling}

\subsection{Short-window cutoff families}
\label{subsec:G.window}

Let $\chi\in \mathcal{S}(\R)$ be even with $\chi(0)=1$.  For $\eta>0$ define the rescaled cutoff
\begin{equation}\label{eq:G.rescaled.cutoff}
\chi_\eta(u)=\chi(u/\eta).
\end{equation}
Its Fourier transform is
\begin{equation}\label{eq:G.rescaled.fourier}
\widehat{\chi_\eta}(t)
=
\eta\,\widehat{\chi}(\eta t).
\end{equation}
This scaling identity is used throughout the microlocal parametrix arguments and in the length-spectrum truncations,
where the effective time support is of size $|t|\ll \eta^{-1}$.

\subsection{Oscillatory test families}
\label{subsec:G.osc}

For $\tau\in\R$ and $\Delta>0$ define
\begin{equation}\label{eq:G.osc.family}
h_{\tau,\Delta}(t)=W(t/\Delta)\,e^{-i\tau t},
\end{equation}
where $W\in C_c^\infty(\R)$ is even.
Then
\begin{equation}\label{eq:G.osc.transform}
\widehat{h_{\tau,\Delta}}(\xi)
=
\Delta\,\widehat{W}\big(\Delta(\xi+\tau)\big),
\end{equation}
which localizes the spectral side to a window of width $\Delta^{-1}$ centered at $-\tau$.
This is the basic Fourier-analytic mechanism behind the trace-to-determinant interface.

\section{A consistency table of factors}
\label{sec:G.table}

For ease of reference we list the key factor conventions that must be synchronized across the monograph.

\begin{center}
\begin{tabular}{ll}
\hline
\textbf{Object} & \textbf{Normalization} \\
\hline
Fourier transform on $\R$ & $\widehat{f}(\xi)=\int_{\R} f(t)e^{-i\xi t}\,dt$ \\
Inverse transform & $f(t)=\frac{1}{2\pi}\int_{\R}\widehat{f}(\xi)e^{i\xi t}\,d\xi$ \\
Plancherel & $\|f\|_2^2=\frac{1}{2\pi}\|\widehat{f}\|_2^2$ \\
Convolution square & $\widehat{f*f^\ast}=|\widehat{f}|^2$ \\
Eigenvalue parameter & $\lambda_j=\frac14+r_j^2$ \\
Continuous spectrum measure & $\frac{1}{4\pi}\sum_{\mathfrak{a}}\int_{\R}(\cdot)\,dr$ \\
Scattering input & $\varphi_{\mathfrak{a}}(1/2+ir)=\frac{d}{ds}\log\varphi_{\mathfrak{a}}(s)$ \\
Window scaling & $\chi_\eta(u)=\chi(u/\eta)$,\quad $\widehat{\chi_\eta}(t)=\eta\widehat{\chi}(\eta t)$ \\
\hline
\end{tabular}
\end{center}

\section{Final remarks}
\label{sec:G.final}

All later identities in this monograph are to be interpreted relative to the conventions fixed in this appendix.
Whenever a formula is imported from the literature, it is implicitly converted to the present normalization by
adjusting Fourier factors and spectral measures.
This ensures that the constants appearing in the main theorems (local Weyl laws, variance bounds, and determinant
interfaces) are stable and reproducible.

\section*{Bibliographic notes}

The Fourier convention adopted here is standard in microlocal analysis.
For the hyperbolic spectral parametrization and scattering normalization, we follow the conventions of
Hejhal~\cite{Hejhal1983} and Iwaniec~\cite{Iwaniec2002}.
For the Selberg--Harish-Chandra transform and the Plancherel measure, see also Buser~\cite{Buser1992}.

% ======================================================================
% End of G-fourier-normalizations-and-measures.tex
% ======================================================================
